% !TEX root = ./main.tex

\chapter{绪论}

\section{研究背景与意义}

近年来，生成式人工智能在图像生成、视频生成、数字人驱动和智能创作等任务中快速发展。生成对抗网络通过对抗训练推动了人脸生成、图像翻译和风格迁移等应用，扩散模型进一步提升了生成质量、文本对齐能力和可控编辑能力，多模态大模型则使图像、视频与自然语言之间的联合理解和生成更加便捷\cite{ref1,ref4,ref6}。这些技术显著降低了数字内容生产门槛，使普通用户也可以借助文本提示、参考图像或简单参数快速生成具有较高真实感的图像和视频。

AIGC技术的普及为数字创意、教育展示、影视制作、广告设计和人机交互带来了新的工具形态，也改变了内容生产和传播方式。相比传统人工制作流程，生成式模型能够在较短时间内完成素材构造、风格变化、局部编辑和动态内容生成，提升了创作效率和表达丰富度。对于科研和工程实践而言，AIGC也提供了构造实验样本、验证检测模型和模拟复杂场景的重要手段。

然而，生成能力提升也使AI生成内容与真实内容之间的边界逐渐模糊。面向人脸的Deepfake、人脸替换、口型同步和数字人驱动技术可能被用于身份冒用、诈骗、舆论操纵和名誉损害；面向通用场景的文生图、图生视频和多模态编辑技术也能够生成新闻事件、灾害画面、公共人物影像或其他具有误导性的视觉内容。一旦这些内容经过社交平台二次压缩、剪辑转发和文字叙事包装，普通用户仅凭直观观察往往难以及时辨别其来源和真实性。

除真实性风险外，AIGC还带来了内容安全风险。生成模型可能在不当提示、越狱诱导或训练数据偏差影响下输出暴力、色情、仇恨、赌博、毒品和敏感符号等违规内容，视频生成还可能通过连续画面形成更强的行为暗示和传播冲击。对平台治理而言，风险并不只来自单张图片或单段视频本身，还来自生成、上传、传播、复用和再编辑构成的完整链路。单纯依赖人工审核难以长期应对海量多模态内容，也难以在不同任务之间保持一致的判定标准。

因此，围绕AIGC内容开展自动化检测与安全分析具有重要研究意义。一方面，虚假内容检测需要识别图像或视频是否由AI生成、是否经过深度伪造或编辑，并结合纹理、频域、语义和时序等线索给出可信判断\cite{ref9,ref10,ref22}。另一方面，不安全内容审核需要关注内容是否包含暴力、色情、仇恨、敏感符号等风险，并输出标签、风险等级、处置建议和原因说明。二者关注点不同，但在实际平台中往往面对相同的图像和视频输入，需要统一组织、统一展示和统一记录。

从工程应用角度看，构建自动化、可扩展的AIGC安全分析工具，有助于把分散的生成能力、检测模型、审核接口和结果管理流程整合为可运行平台。该平台既可以服务于毕业设计中的算法验证和系统实现，也可以为新闻核验、内容审核、教学演示和网络治理提供原型支撑。本文正是在上述背景下，面向图像与视频场景设计并实现多模态内容检测与分析平台，重点解决多源能力接入、真实性检测、安全审核、结果可视化和历史记录管理等问题。

\section{国内外研究现状}

\subsection{AIGC生成技术研究现状}

图像生成技术经历了从生成对抗网络到扩散模型的转变。GAN类方法通过生成器与判别器的对抗训练建模图像分布，在人脸生成、图像翻译和风格迁移中较早取得突破\cite{ref1,ref2,ref3}。扩散模型则通过逐步加噪与逆向去噪生成图像，训练稳定性和生成质量更高，以Stable Diffusion为代表的模型推动了文生图、图像编辑和多条件控制等应用发展\cite{ref4,ref5,ref6,ref7,ref19}。模型能力提升和工具链开源使伪造内容构造门槛进一步降低。

视频生成技术在时序建模、多帧一致性和跨模态对齐方面持续演进，文本到视频、图像到视频和数字人驱动等路线已能生成具备动作和场景变化的短视频。Deepfake、人脸替换、人脸动画和属性编辑技术也在视觉真实感与动作自然度上不断提升，使虚假视频更具欺骗性，并使多模态内容安全问题更加复杂。

\subsection{虚假内容检测研究现状}

AI生成内容检测已形成多条路线。早期方法主要从纹理分布、频域伪影、颜色统计和成像噪声等低层特征出发，利用真实内容与生成内容之间的分布差异进行判别\cite{ref9,ref22}。这类方法解释性较强，但面对新型生成模型时容易泛化不足。

随着深度学习发展，基于卷积神经网络和Transformer的监督检测方法逐渐成为主流\cite{ref13,ref14}。它们能从标注样本中学习复杂判别特征，但对训练数据分布依赖较强。为提升跨模型泛化能力，研究者开始利用大规模预训练视觉模型或视觉语言模型提取通用表征，并通过轻量分类头、近邻检索或重建对比机制完成检测\cite{ref8,ref10,ref12}。在视频方向，逐帧检测、时序建模、局部与全局联合分析以及训练自由检测方法均被用于捕捉帧间不连续、口型不一致和局部结构异常等特征。

\subsection{不安全内容审核研究现状}

不安全内容审核关注内容本身是否涉及暴力、色情、仇恨、敏感符号或其他违规风险。常见方法通常先定义风险标签体系，再通过多标签分类、目标检测、OCR和时序分析提取证据，例如武器、血迹、攻击动作、裸露区域、敏感标志和违规文字等，随后结合阈值、规则或多模态语义推理给出通过、复审或阻断建议。多模态大模型则能够结合视觉输入与自然语言规则，输出风险标签、建议处置和原因说明，适合复杂开放场景。

从部署角度看，单一方法难以同时满足识别范围、推理效率、成本和可维护性要求。本地模型具备自主可控优势，但部署和更新成本较高；外部接口接入便捷、覆盖面较广，却存在依赖和成本问题。因此，如何将本地模型、外部接口和多模态大模型统一组织为可用平台，是当前研究与工程实践中的重要问题。

\subsection{现有研究的不足与问题分析}

综合现有研究，AIGC安全分析仍存在三点不足：图像检测、视频检测和内容审核常分散在不同任务框架中，缺少一体化平台；算法效果与工程可用性存在落差，本地模型、云端接口和多模态大模型缺少统一接入与比较机制；检测结果的解释、风险分级和记录管理仍不完善，难以直接服务教学、科研和平台治理。本文因此从系统集成与工程实现出发，构建面向AIGC安全的多模态内容检测与分析平台。

\section{研究内容与技术路线}

\subsection{主要研究内容}

围绕AIGC图像与视频内容的真实性风险和安全性风险，本文主要开展四方面工作。

第一，完成面向平台实现的需求分析与功能拆解。围绕图像、视频两类输入，梳理虚假内容检测、不安全内容审核、生成样本管理和平台治理等业务场景，明确内容上传、URL输入、模型选择、风险分级、结果展示、历史追踪和一键送检等功能需求，为后续系统架构和接口实现提供约束。

第二，设计并实现多模态内容检测与分析平台。平台采用前端展示层、后端代理层、模型推理与外部能力层分离的结构，实现内容上传、URL输入、图像/视频虚假检测、图像/视频安全审核、内容生成、结果可视化、生成样本管理、检测记录追踪和生成结果一键送检等功能。

第三，研究多源检测能力的组织方式。图像虚假检测接入本地UniversalFakeDetect模型作为可复验基线，开放场景图像/视频检测和安全审核接入多模态接口，并通过统一输入参数、结果字段和风险等级消解不同能力之间的协议差异。

第四，围绕平台可用性与检测链路有效性开展验证。实验结合ProGAN测试样本、平台生成样本和典型图像/视频素材，重点说明当前平台中可稳定运行的能力、适合作为研究基线的能力，以及仍需后续补强的外部依赖能力。

\subsection{技术路线}

本文遵循“问题分析、方法整合、平台实现、实验验证”的技术路线。首先调研AIGC生成与检测方法，明确虚假内容检测和不安全内容识别两类任务；其次结合开源模型与外部审核接口，构建统一检测能力接入框架；再次通过后端代理实现鉴权封装、接口转发、异步轮询和结果标准化处理，并在前端完成上传、检测、展示和管理等交互；最后围绕检测效果、响应特征和平台可用性开展实验分析。

\section{论文特色与创新点}

本文的特色主要体现在三方面。第一，构建多模态内容生成、虚假内容检测与安全审核一体化的AIGC内容安全分析流程，面向图像与视频场景贯通生成样本、送检任务、风险判断和结果展示，使平台能够同时支撑内容生产验证、AI生成鉴别和违规风险识别。第二，实现本地模型与外部接口结合的混合检测架构，用UniversalFakeDetect形成可复验图像检测基线，用多模态接口补充视频理解、开放场景鉴别和安全审核能力。第三，将生成样本管理、一键送检、风险分级、结果可视化和检测历史追踪纳入同一流程，使平台不只返回模型原始响应，而是输出可理解、可展示、可复核的安全分析结果。

需要说明的是，本文创新重点不在于提出全新的底层检测算法，而在于将已有模型、生成能力和审核能力组织为一个可运行、可验证、可扩展的安全分析平台。

\section{论文结构安排}

本文共六章。第一章介绍研究背景、现状、研究内容、技术路线和创新点；第二章梳理深度伪造与多模态内容生成方法、虚假检测、不安全内容审核与判定和系统开发技术；第三章重点分析系统目标、需求、总体架构、业务流程和接口设计；第四章说明多源检测、图像/视频检测、安全审核、生成管理和部署实现；第五章给出实验设计、样本构成、检测验证和可用性分析；第六章总结全文，分析不足并提出后续方向。
